\clearpage
\section{Resource permutation of literal distributors}
\label{sec:copying-distributor-permutation}
\termdef{copying-distributor-permutation}
Graph-level environment renaming must respect the actual distributors of
Figure~16. We prove their resource-permutation law, without deleting their
weakening or contraction edges. This is a prerequisite for bijective renaming
of entire translated terms. It is not the whole-term renaming theorem, and
it does not handle weakening or identification of distinct resources.

\subsection{Permuting resources while retaining use order}
Let $\pi$ be a bijection of the $k$ resource positions. For any component
carrier $X$, define a bijection on tagged elements by
$(i,x)\mapsto(\pi(i),x)$, with inverse using $\pi^{-1}$.
The inverse identities follow directly from those of $\pi$; the component
vertex or edge is unchanged.
\lean{Copying/DistributorPermutation}{9}{14}{Distributor.resourceTags}

Recall that $D_{k,m}$ places leaf $j$ of resource $i$ at input position
$i+kj$ and its output at position $i$. Define the permuted boundary by
precomposing the original input and output bijections with resource retagging.
The resulting input position is $\pi(i)+kj$ and the output position is
$\pi(i)$. Thus each use retains its position within the use-major ordering,
and the same permutation is applied at every use. The original boundary's
typing proofs supply the new boundary's typing proofs. No resource is removed,
identified or duplicated by this operation.
\lean{Copying/DistributorPermutation}{17}{24}{Distributor.permutedBoundary}
\lean{Copying/DistributorPermutation}{27}{31}{Distributor.permuted_input_position}
\lean{Copying/DistributorPermutation}{32}{34}{Distributor.permuted_output_position}

Assemble the original family of canonical contraction trees with this
boundary. Every component retains its internal vertices, labels and ordered
incidences. The original finite enumeration of resource positions is used.
This defines the boundary-permuted forest, including $k=0$ and $m=0$ without
separate non-emptiness assumptions.
\lean{Copying/DistributorPermutation}{37}{41}{Distributor.permutedNet}

\subsection{An isomorphism of the actual finite graphs}
Retagging an internal vertex after embedding it into the permuted forest
has the same result as embedding the unchanged component vertex at resource
$\pi(i)$ in the original forest. For internal vertices this is the definition
of retagging. For inputs and outputs it is the defining boundary permutation.
Importantly, the statement fixes the resulting external port itself; it is
stronger than equality of its wire type or equality of boundary sizes.
\lean{Copying/DistributorPermutation}{44}{49}{Distributor.permuted_embed}

Use resource retagging for both the internal-vertex and edge bijections.
Labels agree because all resource components are copies of the same canonical
tree. For each source list, compose its two maps and apply the preceding
embedding identity to each entry, preserving list order. The same argument
applies to the target list. This constructs an ordinary graph isomorphism
from the permuted forest to the original $D_{k,m}$.
The zero-use case still contains every canonical weakening edge; the
single-use case is not simplified to a wire.
\lean{Copying/DistributorPermutation}{53}{68}{Distributor.permuted_iso}

Relabel the permuted forest into the copying alphabet at any allowance $d$,
using the same primitive embedding as the original distributor. Its carriers
and incidence lists are unchanged. The preceding bijections therefore give
a recursive isomorphism: corresponding embedded primitive labels satisfy
reflexivity of the recursive label relation, and the same source and target
proofs apply. The result holds at every finite box allowance.
\lean{Copying/DistributorPermutation}{71}{75}{Distributor.permutedAtDepth}
\lean{Copying/DistributorPermutation}{77}{85}{Distributor.permuted_deepIso}

\subsection{Separate variable and atom environments}
Apply independent bijections $\rho$ and $\alpha$ to the term and atom forests.
Their ordered tensor is the permuted environment distributor. Term resources
remain before atom resources, and no permutation mixes the two wire sorts.
Tensor congruence assembles the two proved recursive isomorphisms into one
for the original environment distributor.
\lean{Copying/DistributorPermutation}{94}{98}{Translation.permutedEnvironmentDistributor}
\lean{Copying/DistributorPermutation}{101}{105}{Translation.environmentDistributor_permutation}

Finally, take one use and transport along the same interface equalities
used in the original single-use constructor. Isomorphism preservation under
that transport gives the single-use permutation law. All original distributor
edges remain in both graphs.
\lean{Copying/DistributorPermutation}{108}{112}{Translation.singleUse_permutation}

The remaining translation work must propagate boundary renaming through
selected-resource clauses, ordered pairing, binders and their boxes, and all
term constructors. Non-bijective renamings and substitution require additional
sharing arguments. No observational or template theorem follows merely from
this distributor isomorphism; Proposition~12.2 remains unfinished.
